% =====================================================================
% Bab III — Dasar Teori
% =====================================================================
\chapter{Dasar Teori}
\label{bab:dasar-teori}

Bab ini menyajikan fondasi teoretis dan matematis dari metode-metode yang
digunakan dalam penelitian ini. Subbab~\ref{sec:bab3_formulasi_rul} memformulasikan
masalah prediksi RUL dan \emph{health indicator} secara formal.
Subbab~\ref{sec:bab3_bantalan_bpfx} merangkum teori \emph{rolling element bearing} dan
frekuensi karakteristik yang menjadi landasan analisis interpretabilitas.
Subbab~\ref{sec:bab3_analisis_sinyal} membahas teknik analisis sinyal getaran
yang relevan. Subbab~\ref{sec:bab3_lstm_xlstm} menguraikan derivasi matematis
LSTM, sLSTM, dan mLSTM. Subbab~\ref{sec:bab3_ssm_mamba} menjelaskan
\emph{state space model} kontinu dan proses diskritisasi hingga mekanisme
selektivitas Mamba. Subbab~\ref{sec:bab3_sae_interpretabilitas} mendefinisikan
\emph{Sparse Autoencoder}, SHAP, dan \emph{Integrated Gradients}.
Subbab~\ref{sec:bab3_metrik} mendefinisikan metrik evaluasi prognostik
yang digunakan dalam eksperimen.

% -----------------------------------------------------------------
\section{Formulasi Masalah Prediksi RUL dan \emph{Health Indicator}}
\label{sec:bab3_formulasi_rul}
% -----------------------------------------------------------------

Prediksi RUL merupakan tugas regresi yang bertujuan memperkirakan sisa umur
operasi komponen pada setiap titik waktu selama siklus hidupnya.
Subbab ini mendefinisikan secara formal masalah prediksi RUL, skema pelabelan
target, dan konsep \emph{health indicator} yang digunakan sebagai representasi
kondisi \emph{bearing}.

\subsection{Definisi Formal Prediksi RUL}
\label{subsec:bab3_definisi_rul}

Diberikan sekuens fitur kondisi $\mathbf{x}_{1:t} = (\mathbf{x}_1, \mathbf{x}_2,
\ldots, \mathbf{x}_t)$ di mana $\mathbf{x}_\tau \in \mathbb{R}^F$ merupakan
vektor $F$ fitur pada waktu $\tau$, masalah prediksi RUL dirumuskan sebagai
pencarian fungsi $f_\theta$ yang memetakan:

\begin{equation}\label{eq:bab3_rul_mapping}
  \hat{r}_t = f_\theta(\mathbf{x}_{t-L+1:t}),
\end{equation}

\noindent di mana $\hat{r}_t \in [0, 1]$ adalah prediksi RUL ternormalisasi
pada waktu $t$, $L$ adalah panjang \emph{window} input, dan $\theta$ adalah
parameter yang dapat dilatih. Nilai $\hat{r}_t = 1$ menyatakan bahwa komponen
berada dalam kondisi sehat sempurna, sedangkan $\hat{r}_t = 0$ menyatakan
kegagalan (\emph{end of life}).

Target pelatihan $r_t$ didefinisikan menggunakan dua skema pelabelan yang
dipertimbangkan dalam penelitian ini.

\subsection{Skema Pelabelan RUL}
\label{subsec:bab3_skema_label}

\textbf{Skema linier} menurunkan target RUL secara linier dari 1 pada
rekaman pertama hingga 0 pada rekaman terakhir:

\begin{equation}\label{eq:bab3_label_linier}
  r_t = \frac{T - t}{T - 1}, \quad t = 1, 2, \ldots, T,
\end{equation}

\noindent di mana $T$ adalah total jumlah rekaman dalam siklus hidup \emph{bearing}.
Skema ini sederhana namun memiliki kelemahan karena memaksa model belajar
degradasi yang terjadi sejak awal masa pakai, padahal secara fisik \emph{bearing}
berada dalam kondisi sehat untuk periode yang panjang sebelum tanda-tanda
degradasi muncul.

\textbf{Skema \emph{piecewise}} \citetitb{Rui2026xLSTMTransformer}
mempertahankan nilai target konstan pada 1 hingga titik \emph{degradation
onset} $t_d$, kemudian menurun secara linier ke 0 pada \emph{end of life} $t_f$:

\begin{equation}\label{eq:bab3_label_piecewise}
  r_t = \begin{cases}
    1, & t < t_d \\
    \dfrac{t_f - t}{t_f - t_d}, & t_d \leq t \leq t_f
  \end{cases}
\end{equation}

\noindent Titik $t_d$ ditentukan secara otomatis melalui deteksi perubahan
signifikan pada \emph{health indicator} menggunakan metode $3\sigma$ pada
jendela \emph{sliding}. Skema ini lebih realistis secara fisik karena mencerminkan
fase sehat panjang diikuti fase degradasi yang relatif singkat. Pada penelitian
ini, skema \emph{piecewise} dipilih sebagai skema utama.

\subsection{\emph{Health Indicator} dan Ekstraksi Fitur}
\label{subsec:bab3_health_indicator}

\emph{Health indicator} (HI) merupakan representasi ringkas dari kondisi \emph{bearing}
yang diekstraksi dari sinyal getaran mentah. Pada penelitian ini, HI dikonstruksi
dari fitur statistik dan frekuensial yang dihitung pada setiap rekaman. Setelah
ekstraksi, HI diperhalus menggunakan pemulusan eksponensial dengan parameter
$\alpha$:

\begin{equation}\label{eq:bab3_ema}
  \tilde{x}_\tau = \alpha \cdot x_\tau + (1 - \alpha) \cdot \tilde{x}_{\tau-1},
  \quad \alpha \in (0, 1],
\end{equation}

\noindent di mana $x_\tau$ adalah nilai fitur mentah dan $\tilde{x}_\tau$ adalah
nilai yang diperhalus. Pemulusan ini mengurangi \emph{noise} impulsif dalam
sinyal getaran sambil mempertahankan tren degradasi jangka panjang.

Tabel~\ref{tab:bab3_notasi} merangkum seluruh simbol matematika utama yang
digunakan lintas subbab dalam Bab~\ref{bab:dasar-teori} dan
Bab~\ref{bab:metodologi}. Tabel ini berfungsi sebagai referensi cepat untuk
pembaca yang memerlukan klarifikasi notasi saat membaca derivasi formal pada
subbab-subbab berikutnya.

\begin{table}[ht]
  \centering
  \caption{Daftar simbol matematika utama yang digunakan pada Bab~III
           dan Bab~IV beserta dimensi dan keterangannya.}
  \label{tab:bab3_notasi}
  \small
  \begin{tabular}{@{}cll@{}}
    \toprule
    \textbf{Simbol} & \textbf{Dimensi} & \textbf{Keterangan} \\
    \midrule
    $T$              & skalar          & Total rekaman dalam satu siklus hidup \emph{bearing} \\
    $L$              & skalar          & Panjang \emph{window} input sekuens \\
    $F$              & skalar          & Jumlah fitur input (dimensi vektor $\mathbf{x}_t$) \\
    $t$              & skalar          & Indeks waktu (rekaman ke-$t$) \\
    $r_t$            & $[0,1]$         & Label RUL ternormalisasi pada waktu $t$ \\
    $\hat{r}_t$      & $[0,1]$         & Prediksi RUL oleh model pada waktu $t$ \\
    $\mathbf{x}_t$   & $\mathbb{R}^F$  & Vektor fitur input pada waktu $t$ \\
    $\mathbf{h}_t$   & $\mathbb{R}^d$  & \emph{Hidden state} rekuren pada waktu $t$ \\
    $\mathbf{c}_t$   & $\mathbb{R}^d$  & \emph{Cell state} LSTM (vektor) pada waktu $t$ \\
    $\mathbf{C}_t$   & $\mathbb{R}^{d\times d}$ & Memori matriks mLSTM pada waktu $t$ \\
    $\mathbf{A}$     & $\mathbb{R}^{N\times N}$ & Matriks transisi SSM kontinu \\
    $\mathbf{B}$     & $\mathbb{R}^{N\times 1}$ & Matriks proyeksi input SSM \\
    $\mathbf{C}$     & $\mathbb{R}^{1\times N}$ & Matriks proyeksi output SSM \\
    $\Delta_t$       & skalar          & Langkah diskritisasi adaptif pada SSM selektif \\
    $k$              & skalar          & Jumlah unit aktif pada aktivasi Top-$k$ SAE \\
    $d_{\rm lat}$    & skalar          & Dimensi ruang laten ($m$ pada SAE, $d$ pada model) \\
    \bottomrule
  \end{tabular}
\end{table}

% -----------------------------------------------------------------
\section{\emph{Rolling Element Bearing} dan Frekuensi Karakteristik}
\label{sec:bab3_bantalan_bpfx}
% -----------------------------------------------------------------

\emph{Rolling element bearing} merupakan
komponen mekanik yang memungkinkan rotasi poros dengan gesekan minimal. Pada
\emph{bearing} tipe bola alur dalam (\emph{deep groove ball bearing}), terdapat
empat komponen utama: cincin luar (\emph{outer race}), cincin dalam
(\emph{inner race}), \emph{rolling elements}, dan sangkar
(\emph{cage}). Kegagalan \emph{bearing} umumnya dimulai sebagai spalling pada salah
satu komponen, yang menghasilkan impuls periodik dalam sinyal getaran
\citetitb{Lei2018Machinery}. Gambar~\ref{fig:bab3_bearing_geom} mengilustrasikan
penampang lintang \emph{bearing} bola alur dalam dengan parameter geometri yang
menjadi dasar perhitungan frekuensi karakteristik.

\begin{figure}[ht]
  \centering
  \begin{tikzpicture}[scale=1.0,
      thick,
      raceline/.style={draw, very thick},
      dim/.style={draw, thin, <->, >=Stealth, font=\scriptsize},
      lbl/.style={font=\small},
    ]
    % Outer race ring
    \draw[raceline, fill=gray!20] (0,0) circle (2.8cm);
    \draw[raceline, fill=white]   (0,0) circle (2.2cm);
    % Inner race ring
    \draw[raceline, fill=gray!20] (0,0) circle (1.6cm);
    \draw[raceline, fill=white]   (0,0) circle (0.9cm);
    % Rolling elements (balls) — 8 at equal angles
    \foreach \ang in {0, 45, 90, 135, 180, 225, 270, 315}{
      \draw[fill=gray!40] (\ang:1.9cm) circle (0.28cm);
    }
    % Cage lines (dashed circle at pitch diameter)
    \draw[dashed, thin] (0,0) circle (1.9cm);
    % Dimension D (pitch diameter)
    \draw[dim] (0,0) -- node[above, pos=0.6]{\footnotesize $D/2$} (0:1.9cm);
    % Dimension d (rolling element diameter)
    \draw[dim] (0:1.9cm) -- node[above, font=\scriptsize]{$d$} (0:2.18cm);
    % Labels
    \node[lbl, above=0.15cm] at (0:2.5cm)   {\footnotesize Cincin luar};
    \node[lbl, above=0.15cm] at (180:1.9cm) {\footnotesize Sangkar};
    \node[lbl, below=0.2cm]  at (180:1.25cm){\footnotesize Cincin dalam};
    \node[lbl, font=\scriptsize] at (90:1.9cm) {bola ($n=8$)};
  \end{tikzpicture}
  \caption{Penampang lintang \emph{bearing} bola alur dalam (\emph{deep groove ball
           bearing}) dengan anotasi parameter geometri: $D$ = diameter lingkaran
           \emph{pitch}, $d$ = diameter \emph{rolling element}, $n$ = jumlah \emph{rolling element}. Nilai $D$, $d$, $n$, dan sudut kontak $\theta$ digunakan
           untuk menghitung BPFO, BPFI, BSF, dan FTF
           (Persamaan~\ref{eq:bab3_bpfo}--\ref{eq:bab3_ftf}).}
  \label{fig:bab3_bearing_geom}
\end{figure}

\subsection{Frekuensi Karakteristik \emph{Bearing}}
\label{subsec:bab3_bpfx}

Saat terjadi spalling pada komponen \emph{bearing}, impuls periodik yang ditimbulkan
memiliki frekuensi yang dapat diprediksi dari geometri \emph{bearing} dan kecepatan
rotasi poros. Empat frekuensi karakteristik utama didefinisikan sebagai
berikut \citetitb{Lei2018Machinery}.

Misalkan $f_r$ = frekuensi rotasi poros (Hz), $n$ = jumlah \emph{rolling element},
$d$ = diameter \emph{rolling element} (mm), $D$ = diameter lingkaran \emph{pitch}
(mm), dan $\theta$ = sudut kontak (rad). Frekuensi lintasan bola pada cincin
luar (\emph{Ball Pass Frequency Outer race}, BPFO) adalah:

\begin{equation}\label{eq:bab3_bpfo}
  \text{BPFO} = \frac{n}{2} f_r \left(1 - \frac{d}{D}\cos\theta\right).
\end{equation}

Frekuensi lintasan bola pada cincin dalam (\emph{Ball Pass Frequency Inner
race}, BPFI) adalah:

\begin{equation}\label{eq:bab3_bpfi}
  \text{BPFI} = \frac{n}{2} f_r \left(1 + \frac{d}{D}\cos\theta\right).
\end{equation}

Frekuensi putar \emph{rolling element} (\emph{Ball Spin Frequency}, BSF) adalah:

\begin{equation}\label{eq:bab3_bsf}
  \text{BSF} = \frac{D}{2d} f_r \left[1 - \left(\frac{d}{D}\cos\theta\right)^2\right].
\end{equation}

Frekuensi fundamental sangkar (\emph{Fundamental Train Frequency} atau
\emph{Cage Frequency}, FTF) adalah:

\begin{equation}\label{eq:bab3_ftf}
  \text{FTF} = \frac{1}{2} f_r \left(1 - \frac{d}{D}\cos\theta\right).
\end{equation}

Pada \emph{bearing} \emph{deep groove ball bearing} dengan sudut kontak yang kecil,
asumsi $\theta \approx 0$ ($\cos\theta \approx 1$) memberikan penyederhanaan
yang cukup akurat untuk perhitungan estimasi.

\subsection{Geometri \emph{Bearing} Benchmark}
\label{subsec:bab3_geometri_benchmark}

Tabel~\ref{tab:bab3_geometri_bearing} menyajikan parameter geometri dan
frekuensi karakteristik teoretis untuk dua \emph{bearing} yang digunakan pada
dataset benchmark penelitian ini.

\begin{table}[ht]
  \centering
  \caption{Geometri \emph{bearing} benchmark dan frekuensi karakteristik teoretis
           (asumsi $\theta \approx 0$).}
  \label{tab:bab3_geometri_bearing}
  \small
  \begin{tabular}{@{}lccccccc@{}}
    \toprule
    \emph{Bearing} & Dataset & $n$ & $d$ (mm) & $D$ (mm) & $f_r$ (Hz)
      & BPFO (Hz) & BPFI (Hz) \\
    \midrule
    NSK 6804 & PHM2012 & 13 & 3{,}50 & 25{,}50 & 30{,}0 & 168{,}3 & 228{,}7 \\
    LDK UER204 & XJTU-SY & 8 & 7{,}92 & 34{,}55 & 35{,}0 & 107{,}9 & 172{,}1 \\
    \bottomrule
  \end{tabular}

  \smallskip
  \footnotesize
  \emph{Catatan}: Nilai frekuensi di atas merupakan estimasi berdasarkan data
  geometri yang tersedia; verifikasi dengan datasheet aktual direkomendasikan
  sebelum penggunaan diagnosis klinis.
  % TODO(verify-with-pembimbing): konfirmasi geometri NSK 6804 dari datasheet
  % resmi; beberapa sumber berbeda pada nilai d dan D untuk model ini.
\end{table}

Frekuensi-frekuensi karakteristik pada Tabel~\ref{tab:bab3_geometri_bearing}
digunakan dalam analisis interpretabilitas Bab~\ref{bab:hasil-pembahasan}
sebagai referensi untuk mengevaluasi apakah neuron SAE yang paling aktif
berkorespondensi dengan frekuensi-frekuensi yang relevan secara fisik.

% -----------------------------------------------------------------
\section{Analisis Sinyal Getaran}
\label{sec:bab3_analisis_sinyal}
% -----------------------------------------------------------------

Sinyal getaran \emph{bearing} merupakan superposis dari komponen yang berasal dari
proses fisika yang berbeda: rotasi poros, meshing roda gigi (bila ada), resonansi
struktural, dan impuls periodik akibat cacat \emph{bearing}. Ekstraksi fitur yang
efektif memerlukan pemahaman tentang karakteristik sinyal ini.

\subsection{Fitur Energi Band Frekuensi}
\label{subsec:bab3_band_energy}

Energi RMS (\emph{Root Mean Square}) pada band frekuensi tertentu dihitung
dari spektrum daya sinyal. Untuk rekaman getaran $x[n]$, $n = 0, 1, \ldots, N-1$,
dengan transformasi Fourier $X[k]$, energi pada band $[f_{\rm low}, f_{\rm high}]$
adalah:

\begin{equation}\label{eq:bab3_band_energy}
  E_{\rm band} = \sqrt{\frac{1}{|B|} \sum_{k \in B} |X[k]|^2},
\end{equation}

\noindent di mana $B = \{k : f_{\rm low} \leq k\Delta f \leq f_{\rm high}\}$ dan
$\Delta f = f_s/N$ adalah resolusi frekuensi. Fitur band-energy pada rentang
yang mencakup BPFO, BPFI, dan harmoniknya merupakan indikator yang sensitif
terhadap perkembangan cacat \emph{bearing}.

\subsection{Analisis Spektrum \emph{Envelope} Hilbert}
\label{subsec:bab3_hilbert}

Analisis \emph{envelope} berbasis transformasi Hilbert merupakan teknik standar
untuk mendeteksi modulasi amplitudo pada sinyal getaran impulsif.
Transformasi Hilbert dari sinyal $x(t)$ didefinisikan sebagai:

\begin{equation}\label{eq:bab3_hilbert}
  \mathcal{H}\{x(t)\} = \frac{1}{\pi} \, \text{P.V.} \int_{-\infty}^{\infty}
    \frac{x(\tau)}{t - \tau} \, d\tau,
\end{equation}

\noindent di mana P.V. menunjukkan nilai prinsipal Cauchy. Sinyal analitik
$z(t) = x(t) + j\mathcal{H}\{x(t)\}$ memiliki \emph{envelope} (amplitudo
sesaat) $a(t) = |z(t)| = \sqrt{x(t)^2 + \mathcal{H}\{x(t)\}^2}$. Spektrum
FFT dari $a(t)$ --- yang disebut spektrum \emph{envelope} --- menunjukkan
frekuensi modulasi impuls, yang langsung berkaitan dengan BPFO, BPFI, BSF,
atau FTF jika cacat terlokalisasi pada komponen tertentu \citetitb{Lei2018Machinery}.

% -----------------------------------------------------------------
\section{LSTM dan xLSTM}
\label{sec:bab3_lstm_xlstm}
% -----------------------------------------------------------------

\emph{Long Short-Term Memory} (LSTM) dan variannya xLSTM merupakan fondasi
arsitektur \emph{baseline} dan komponen utama arsitektur usulan dalam
penelitian ini. Gambar~\ref{fig:bab3_evolusi_seq} merangkum silsilah evolusi
dua keluarga arsitektur sekuens yang berperan dalam disertasi ini: keluarga
rekuren berbasis LSTM dan keluarga \emph{state-space model} (SSM) yang berkulminasi
pada Mamba. Pemahaman terhadap jalur evolusi ini penting untuk
mengidentifikasi limitasi tiap generasi dan justifikasi desain hibrida
yang diuraikan pada Bab~\ref{bab:metodologi}.

\begin{figure}[ht]
  \centering
  \resizebox{\textwidth}{!}{%
  \begin{tikzpicture}[
      node distance = 0.55cm and 1.4cm,
      evo/.style  = {rectangle, rounded corners=3pt, draw, thick,
                     fill=white, align=center, minimum height=0.75cm,
                     minimum width=2.6cm, font=\small},
      final/.style = {evo, fill=blue!12, very thick},
      arr/.style  = {-{Stealth[length=5pt]}, thick},
      lane/.style = {draw=gray!40, rounded corners, dashed, inner sep=4pt},
    ]
    \node[evo]   (rnn)   {RNN (1986)};
    \node[evo,   below=of rnn]  (lstm)  {LSTM (1997)};
    \node[evo,   below=of lstm] (bi)    {BiLSTM / GRU};
    \node[final, below=of bi]   (xlstm) {xLSTM (2024)\\{\scriptsize sLSTM + mLSTM}};

    \draw[arr] (rnn)  -- (lstm);
    \draw[arr] (lstm) -- (bi);
    \draw[arr] (bi)   -- (xlstm);

    \begin{pgfonlayer}{background}
      \node[lane, fit=(rnn)(xlstm), label=above:{\footnotesize Keluarga Rekuren}] {};
    \end{pgfonlayer}

    \node[evo,   right=3.2cm of rnn]   (hippo) {HiPPO / S4 (2020--22)};
    \node[evo,   below=of hippo]        (s4d)   {S4D / DSS (2023)};
    \node[evo,   below=of s4d]          (mamba) {Mamba (2023)};
    \node[final, below=of mamba]        (bimamba){BiMamba (2024)\\{\scriptsize SSM Selektif, Bidir.}};

    \draw[arr] (hippo)  -- (s4d);
    \draw[arr] (s4d)    -- (mamba);
    \draw[arr] (mamba)  -- (bimamba);

    \begin{pgfonlayer}{background}
      \node[lane, fit=(hippo)(bimamba), label=above:{\footnotesize Keluarga SSM}] {};
    \end{pgfonlayer}

    \coordinate (xlstm_bimamba_mid) at ($(xlstm)!0.5!(bimamba)$);
    \node[final, below=1.2cm of xlstm_bimamba_mid] (hybrid)
          {Mamba-xLSTM-Net (Kap.~IV)\\{\scriptsize Arsitektur usulan}};
    \draw[arr, very thick] (xlstm.south)   -- (hybrid.north west);
    \draw[arr, very thick] (bimamba.south) -- (hybrid.north east);
  \end{tikzpicture}%
  }
  \caption{Pohon evolusi arsitektur sekuens yang relevan untuk penelitian ini.
           Keluarga rekuren (kiri) berkembang dari RNN menuju xLSTM dengan memori
           matriks. Keluarga SSM (kanan) berkembang dari HiPPO menuju Mamba
           selektif bidireksional. Mamba-xLSTM-Net (Bab~\ref{bab:metodologi})
           mengintegrasikan kedua cabang ini melalui mekanisme \emph{gated fusion}.}
  \label{fig:bab3_evolusi_seq}
\end{figure}

\subsection{LSTM Klasik}
\label{subsec:bab3_lstm}

LSTM \citetitb{Hochreiter1997LSTM} mengatasi masalah \emph{vanishing gradient}
dalam RNN dengan memperkenalkan \emph{cell state} $\mathbf{c}_t$ yang dilindungi
oleh tiga gerbang multiplikatif. Persamaan LSTM pada waktu $t$ adalah:

\begin{align}
  \mathbf{i}_t &= \sigma(\mathbf{W}_i \mathbf{x}_t + \mathbf{R}_i \mathbf{h}_{t-1}
    + \mathbf{b}_i), \label{eq:bab3_lstm_i} \\
  \mathbf{f}_t &= \sigma(\mathbf{W}_f \mathbf{x}_t + \mathbf{R}_f \mathbf{h}_{t-1}
    + \mathbf{b}_f), \label{eq:bab3_lstm_f} \\
  \mathbf{o}_t &= \sigma(\mathbf{W}_o \mathbf{x}_t + \mathbf{R}_o \mathbf{h}_{t-1}
    + \mathbf{b}_o), \label{eq:bab3_lstm_o} \\
  \tilde{\mathbf{c}}_t &= \tanh(\mathbf{W}_c \mathbf{x}_t + \mathbf{R}_c
    \mathbf{h}_{t-1} + \mathbf{b}_c), \label{eq:bab3_lstm_ctilde} \\
  \mathbf{c}_t &= \mathbf{f}_t \odot \mathbf{c}_{t-1}
    + \mathbf{i}_t \odot \tilde{\mathbf{c}}_t, \label{eq:bab3_lstm_c} \\
  \mathbf{h}_t &= \mathbf{o}_t \odot \tanh(\mathbf{c}_t), \label{eq:bab3_lstm_h}
\end{align}

\noindent di mana $\mathbf{x}_t \in \mathbb{R}^F$ adalah input, $\mathbf{h}_t
\in \mathbb{R}^d$ adalah \emph{hidden state}, $\mathbf{c}_t \in \mathbb{R}^d$
adalah \emph{cell state}, $\sigma$ adalah fungsi sigmoid, $\odot$ adalah
perkalian elemen demi elemen, dan $\mathbf{W}_*, \mathbf{R}_*, \mathbf{b}_*$
adalah parameter yang dapat dilatih. Gerbang $\mathbf{i}_t$, $\mathbf{f}_t$,
$\mathbf{o}_t$ bernilai dalam $[0, 1]$ dan mengontrol aliran informasi ke
dalam, pelupaan, dan keluaran sel memori secara berurutan.

\textbf{Limitasi mekanistik LSTM.} Meskipun LSTM berhasil mengurangi masalah
\emph{vanishing gradient} dibandingkan RNN sederhana, penggunaan fungsi
sigmoid pada ketiga gerbang menghadirkan batasan mekanistik yang penting.
Secara analitis, turunan sigmoid $\sigma'(z) = \sigma(z)(1 - \sigma(z))$
memiliki nilai maksimum $0{,}25$ pada $z = 0$ dan menuju nol secara
eksponensial untuk $|z| \gg 0$. Kondisi \emph{saturation} ini berarti bahwa
ketika terjadi lonjakan amplitudo getaran yang besar dan tiba-tiba menjelang
kegagalan akhir \emph{bearing} --- tepat kondisi ketika model harus merespons dengan
revisi besar terhadap estimasi RUL --- gerbang LSTM yang sudah tersaturasi
menghasilkan gradien yang mendekati nol dan perubahan memori sel yang sangat
kecil \citetitb{Bengio1994Vanishing,Beck2024xLSTM}. Dengan kata lain, LSTM
paling tidak responsif justru pada saat paling dibutuhkan. Batasan ini tidak
dapat diatasi dengan hanya memperdalam arsitektur atau menyetel ulang
hiperparameter, karena sifatnya bersumber dari fungsi aktivasi itu sendiri.
Pengamatan ini memotivasi pengenalan \emph{exponential gating} pada sLSTM,
yang diuraikan pada subbab berikut.

\subsection{sLSTM: \emph{Exponential Gating} dan \emph{Normalizer State}}
\label{subsec:bab3_slstm}

sLSTM \citetitb{Beck2024xLSTM} memperkenalkan dua modifikasi kunci terhadap
LSTM klasik. Pertama, gerbang \emph{input} dan \emph{forget} menggunakan
\emph{exponential activation} sebagai pengganti sigmoid, memungkinkan nilai
gate melampaui rentang $[0, 1]$. Kedua, \emph{normalizer state} $n_t$ ditambahkan
untuk stabilisasi numerik. Persamaan sLSTM pada waktu $t$ adalah:

\begin{align}
  \mathbf{z}_t &= \tanh(\mathbf{W}_z \mathbf{x}_t + \mathbf{R}_z
    \mathbf{h}_{t-1} + \mathbf{b}_z), \label{eq:bab3_slstm_z} \\
  \mathbf{i}_t &= \exp(\mathbf{W}_i \mathbf{x}_t + \mathbf{R}_i
    \mathbf{h}_{t-1} + \mathbf{b}_i), \label{eq:bab3_slstm_i} \\
  \mathbf{f}_t &= \exp(\mathbf{W}_f \mathbf{x}_t + \mathbf{R}_f
    \mathbf{h}_{t-1} + \mathbf{b}_f), \label{eq:bab3_slstm_f} \\
  \mathbf{o}_t &= \sigma(\mathbf{W}_o \mathbf{x}_t + \mathbf{R}_o
    \mathbf{h}_{t-1} + \mathbf{b}_o), \label{eq:bab3_slstm_o} \\
  \mathbf{c}_t &= \mathbf{f}_t \odot \mathbf{c}_{t-1}
    + \mathbf{i}_t \odot \mathbf{z}_t, \label{eq:bab3_slstm_c} \\
  \mathbf{n}_t &= \mathbf{f}_t \odot \mathbf{n}_{t-1}
    + \mathbf{i}_t, \label{eq:bab3_slstm_n} \\
  \mathbf{h}_t &= \mathbf{o}_t \odot \frac{\mathbf{c}_t}{\mathbf{n}_t}.
    \label{eq:bab3_slstm_h}
\end{align}

\emph{Exponential gating} pada Persamaan~\ref{eq:bab3_slstm_i} dan
\ref{eq:bab3_slstm_f} merupakan keunggulan utama sLSTM: model dapat merespons
perubahan tajam pada input dengan merevisi memori sel secara agresif, yang
secara khusus relevan untuk fase akhir degradasi \emph{bearing} di mana amplitudo
getaran meningkat secara dramatis dalam waktu singkat.

\textbf{Justifikasi memori matriks.} Satu dimensi yang belum diatasi oleh
sLSTM adalah kapasitas representasi memori itu sendiri. Memori sel $\mathbf{c}_t
\in \mathbb{R}^d$ merupakan sebuah vektor dengan $d$ entri, yang membatasi
jumlah asosiasi berbeda yang dapat disimpan secara bersamaan menjadi paling
banyak $\mathcal{O}(d)$ entri independen. Dalam konteks prediksi RUL \emph{bearing},
model harus secara bersamaan mengingat pola dari banyak fitur kondisi yang
berbeda (kurtosis, RMS, fitur frekuensi band-energy) dan korelasinya yang
berubah sepanjang siklus degradasi. Vektor $d$-dimensi tunggal tidak cukup
untuk merepresentasikan semua asosiasi fitur-RUL yang relevan secara simultan.
Kapasitas yang terbatas ini menjadi hambatan pada sekuens degradasi yang panjang
dan multi-fitur, di mana model perlu mempertahankan memori tentang tren dari
banyak karakteristik getaran berbeda secara bersamaan. Keterbatasan inilah yang
memotivasi mLSTM dengan memori matriks $\mathbf{C}_t \in \mathbb{R}^{d \times d}$,
yang menyediakan kapasitas $\mathcal{O}(d^2)$ untuk asosiasi kunci-nilai.

\subsection{mLSTM: Memori Matriks}
\label{subsec:bab3_mlstm}

mLSTM \citetitb{Beck2024xLSTM} memperluas kapasitas memori LSTM dari vektor
$\mathbf{c}_t \in \mathbb{R}^d$ ke matriks $\mathbf{C}_t \in \mathbb{R}^{d \times d}$,
terinspirasi oleh memori asosiatif. Setiap \emph{time step} mengasosiasikan
pasangan kunci-nilai $(\mathbf{k}_t, \mathbf{v}_t)$ dalam memori matriks.
Persamaan mLSTM pada waktu $t$ adalah:

\begin{align}
  \mathbf{q}_t &= \mathbf{W}_q \mathbf{x}_t + \mathbf{b}_q, \quad
  \mathbf{k}_t = \mathbf{W}_k \mathbf{x}_t + \mathbf{b}_k, \quad
  \mathbf{v}_t = \mathbf{W}_v \mathbf{x}_t + \mathbf{b}_v,
    \label{eq:bab3_mlstm_qkv} \\
  i_t &= \exp(\mathbf{w}_i^\top \mathbf{x}_t + b_i), \quad
  f_t = \exp(\mathbf{w}_f^\top \mathbf{x}_t + b_f),
    \label{eq:bab3_mlstm_gates} \\
  \mathbf{o}_t &= \sigma(\mathbf{W}_o \mathbf{x}_t + \mathbf{b}_o),
    \label{eq:bab3_mlstm_o} \\
  \mathbf{C}_t &= f_t \cdot \mathbf{C}_{t-1}
    + i_t \cdot \frac{\mathbf{v}_t \mathbf{k}_t^\top}{\sqrt{d}},
    \label{eq:bab3_mlstm_C} \\
  \mathbf{n}_t &= f_t \cdot \mathbf{n}_{t-1} + i_t \cdot \mathbf{k}_t,
    \label{eq:bab3_mlstm_n} \\
  \mathbf{h}_t &= \mathbf{o}_t \odot \frac{\mathbf{C}_t \mathbf{q}_t}{\max(|\mathbf{n}_t^\top \mathbf{q}_t|, 1)}.
    \label{eq:bab3_mlstm_h}
\end{align}

Tidak seperti sLSTM, mLSTM tidak memiliki koneksi rekurens pada gate
(hanya bergantung pada $\mathbf{x}_t$), sehingga komputasinya lebih mudah
diparalelkan. Kapasitas memori matriks $\mathbf{C}_t$ yang berukuran $d \times d$
memberikan kemampuan representasi yang secara kuadratik lebih besar dari
memori vektor LSTM klasik, yang berguna ketika banyak asosiasi fitur vibrasi
perlu dipertahankan sepanjang siklus degradasi panjang.

% -----------------------------------------------------------------
\section{\emph{State Space Model} dan Mamba}
\label{sec:bab3_ssm_mamba}
% -----------------------------------------------------------------

\emph{State space model} (SSM) adalah kerangka matematis untuk pemodelan
sistem dinamis linier yang memiliki akar kuat dalam teori kontrol dan pemrosesan
sinyal. Penerapannya ke \emph{deep learning} membuka peluang untuk pemodelan
sekuens panjang yang efisien.

\subsection{SSM Kontinu}
\label{subsec:bab3_ssm_kontinu}

SSM kontinu-waktu didefinisikan oleh persamaan diferensial orde pertama:

\begin{equation}\label{eq:bab3_ssm_ode}
  \frac{d\mathbf{h}(t)}{dt} = \mathbf{A} \mathbf{h}(t) + \mathbf{B} x(t),
  \qquad y(t) = \mathbf{C} \mathbf{h}(t),
\end{equation}

\noindent di mana $\mathbf{h}(t) \in \mathbb{R}^N$ adalah \emph{state}
tersembunyi, $x(t) \in \mathbb{R}$ adalah sinyal input skalar,
$y(t) \in \mathbb{R}$ adalah output, dan $\mathbf{A} \in \mathbb{R}^{N \times N}$,
$\mathbf{B} \in \mathbb{R}^{N \times 1}$, $\mathbf{C} \in \mathbb{R}^{1 \times N}$
adalah matriks parameter. S4 \citetitb{Gu2022S4} mengusulkan parameterisasi
khusus pada $\mathbf{A}$ (matriks HiPPO) yang memungkinkan SSM merangkum
riwayat input secara teoritis optimal.

\textbf{Limitasi Transformer yang mendorong pengembangan SSM.} Untuk memahami
mengapa SSM menjadi alternatif yang menarik, penting untuk meninjau keterbatasan
arsitektur dominan yang ingin diatasi. Transformer \citetitb{Vaswani2017Attention}
menghitung \emph{self-attention} dengan kompleksitas $\mathcal{O}(L^2 d)$
terhadap panjang sekuens $L$: setiap pasangan posisi dalam sekuens perlu
dibandingkan satu sama lain. Evaluasi empiris yang komprehensif oleh
\citetitb{Tay2022LRA} menunjukkan bahwa meskipun Transformer unggul pada
sekuens pendek dan menengah, performa dan efisiensinya mengalami degradasi
signifikan untuk sekuens panjang ($L > 4.096$). Jejak memori KV-cache tumbuh
linier dengan panjang sekuens dan tidak skalabel untuk aplikasi \emph{streaming}
atau \emph{continual monitoring} seperti pemantauan kondisi mesin berbasis sensor.
Selain itu, mekanisme \emph{masked attention} yang natural untuk pemodelan bahasa
autoregresif tidak secara langsung cocok untuk sinyal kondisi mesin yang dianalisis
pada jendela tetap. SSM dengan kompleksitas $\mathcal{O}(L)$ dan struktur rekurens
yang natural untuk data berbasis waktu memberikan solusi matematika yang lebih
sesuai untuk domain pemantauan kondisi industri.

\subsection{Diskritisasi \emph{Zero-Order Hold}}
\label{subsec:bab3_zoh}

Untuk implementasi pada data diskrit dengan langkah waktu $\Delta$, SSM
kontinu didiskritkan menggunakan metode \emph{zero-order hold} (ZOH):

\begin{align}
  \bar{\mathbf{A}} &= \exp(\Delta \cdot \mathbf{A}),
    \label{eq:bab3_zoh_A} \\
  \bar{\mathbf{B}} &= (\Delta \cdot \mathbf{A})^{-1}
    (\exp(\Delta \cdot \mathbf{A}) - \mathbf{I}) \cdot \Delta \cdot \mathbf{B}.
    \label{eq:bab3_zoh_B}
\end{align}

Model diskrit yang dihasilkan mengikuti rekurens:

\begin{equation}\label{eq:bab3_ssm_diskrit}
  \mathbf{h}_t = \bar{\mathbf{A}} \mathbf{h}_{t-1} + \bar{\mathbf{B}} x_t,
  \qquad y_t = \mathbf{C} \mathbf{h}_t.
\end{equation}

Persamaan~\ref{eq:bab3_ssm_diskrit} dapat dikomputasi baik sebagai rekurens
(untuk inferensi sekuensial) maupun sebagai konvolusi global
$\mathbf{y} = \mathbf{x} * \bar{\mathbf{K}}$ di mana
$\bar{\mathbf{K}} = (\mathbf{C}\bar{\mathbf{B}},
\mathbf{C}\bar{\mathbf{A}}\bar{\mathbf{B}}, \ldots)$ adalah kernel konvolusi
implisit. Dualitas ini memungkinkan pelatihan yang efisien menggunakan konvolusi
(paralel) dan inferensi yang efisien menggunakan rekurens.

\subsection{Mamba: SSM Selektif}
\label{subsec:bab3_mamba_selektif}

Kontribusi utama Mamba \citetitb{Gu2023Mamba} adalah menjadikan parameter
$\Delta$, $\mathbf{B}$, dan $\mathbf{C}$ sebagai fungsi dari input:

\begin{equation}\label{eq:bab3_mamba_selective}
  \mathbf{B}_t = \text{Linear}_B(\mathbf{x}_t), \quad
  \mathbf{C}_t = \text{Linear}_C(\mathbf{x}_t), \quad
  \Delta_t = \text{softplus}(\text{Linear}_\Delta(\mathbf{x}_t)).
\end{equation}

Dengan selektivitas ini, model dapat secara adaptif menentukan informasi
mana yang disimpan dalam \emph{state} tersembunyi dan mana yang diabaikan,
bergantung pada konten input aktual. Nilai $\Delta_t$ yang besar menyebabkan
SSM lebih terfokus pada input saat ini (perilaku mirip konvolusi), sedangkan
$\Delta_t$ yang kecil menyebabkan SSM lebih bergantung pada \emph{state}
sebelumnya (perilaku mirip rekurens). Fleksibilitas ini memungkinkan Mamba
berperilaku adaptif terhadap karakteristik sinyal yang berubah sepanjang
waktu, seperti sinyal degradasi \emph{bearing} yang bersifat non-stasioner.

Mamba-2 \citetitb{Dao2024Mamba2} memperbaiki efisiensi komputasi Mamba
melalui kerangka SSD (\emph{Structured State Space Duality}) yang menunjukkan
kesetaraan antara SSM selektif dan \emph{linear attention}. Pembuktian ini
memungkinkan adopsi teknik optimasi perangkat keras Transformer (seperti
komputasi terstruktur pada GPU) ke dalam SSM, menghasilkan implementasi yang
2--8$\times$ lebih cepat.

\textbf{Mengapa bidireksional diperlukan.} Mamba dalam konfigurasi standar
bersifat unidireksional: rekurens $\mathbf{h}_t = \bar{\mathbf{A}}_t
\mathbf{h}_{t-1} + \bar{\mathbf{B}}_t x_t$ hanya dapat mengakses konteks
historis $\mathbf{x}_{1:t-1}$. Untuk pemodelan bahasa yang bersifat
autoregresif (memprediksi token berikutnya dari token sebelumnya), sifat
kausal ini merupakan properti yang diinginkan. Namun konteks berbeda berlaku
untuk prediksi RUL pada jendela sinyal degradasi: model diberikan \emph{seluruh}
jendela waktu $\mathbf{x}_{t-L+1:t}$ sekaligus dan diminta menghasilkan satu
estimasi RUL --- tidak ada batasan kausalitas yang perlu dipatuhi. Dalam
situasi ini, memproses sekuens hanya dalam satu arah berarti posisi awal
jendela mengakses konteks yang lebih sedikit dibandingkan posisi akhir.
Penelitian pada BERT \citetitb{Devlin2019BERT} dan model bidireksional lainnya
telah membuktikan bahwa konteks dua arah menghasilkan representasi yang
lebih kaya untuk tugas yang bukan autoregresif. Untuk sinyal degradasi
\emph{bearing}, fitur pada jendela awal mengandung informasi tentang kondisi dasar
(\emph{baseline}) yang relevan untuk menginterpretasikan fitur pada jendela
akhir; demikian pula sebaliknya. Oleh karena itu, penelitian ini menggunakan
konfigurasi BiMamba yang memproses jendela dalam dua arah secara paralel,
sebagaimana diuraikan pada subbab berikut.

\subsection{Mamba Bidireksional}
\label{subsec:bab3_bimamba}

Untuk memperkaya representasi dengan konteks dua arah, penelitian ini
menggunakan konfigurasi Mamba bidireksional (BiMamba) di mana dua SSM selektif
memproses sekuens masing-masing dalam arah maju dan mundur:

\begin{align}
  \mathbf{h}^{\rightarrow}_t &= \text{Mamba}_{\rm fwd}(\mathbf{x}_{1:t}),
    \label{eq:bab3_bimamba_fwd} \\
  \mathbf{h}^{\leftarrow}_t &= \text{Mamba}_{\rm bwd}(\mathbf{x}_{t:L}),
    \label{eq:bab3_bimamba_bwd} \\
  \mathbf{h}^{\rm BiMamba}_t &= \text{Linear}([\mathbf{h}^{\rightarrow}_t \;
    ; \; \mathbf{h}^{\leftarrow}_t]),
    \label{eq:bab3_bimamba_fused}
\end{align}

\noindent di mana $[\cdot ; \cdot]$ menunjukkan konkatenasi dan $\text{Linear}$
adalah proyeksi linier dari $2d_{\rm model}$ kembali ke $d_{\rm model}$.
Konfigurasi bidireksional memungkinkan setiap \emph{time step} mengakses
konteks sebelum dan setelahnya, yang berguna untuk memodelkan korelasi
temporal non-kausal dalam \emph{health indicator}.

% -----------------------------------------------------------------
\section{\emph{Sparse Autoencoder} dan Interpretabilitas}
\label{sec:bab3_sae_interpretabilitas}
% -----------------------------------------------------------------

Interpretabilitas mekanistik bertujuan memahami internal komputasi model
\emph{deep learning} pada tingkat representasi tersembunyi, bukan sekadar
mengatribusikan prediksi ke fitur input. Subbab ini mendefinisikan instrumen
utama yang digunakan: SAE, SHAP, dan \emph{Integrated Gradients}.

\subsection{\emph{Sparse Autoencoder} (SAE) dengan Top-\texorpdfstring{$k$}{k}}
\label{subsec:bab3_sae}

Diberikan representasi tersembunyi $\mathbf{z} \in \mathbb{R}^d$ dari
lapisan tertentu model RUL, SAE terdiri dari:

\begin{enumerate}
  \item \emph{Encoder}: $\mathbf{a} = \text{TopK}(\mathbf{W}_{\rm enc}
    \mathbf{z} + \mathbf{b}_{\rm enc}, k)$
  \item \emph{Decoder}: $\hat{\mathbf{z}} = \mathbf{W}_{\rm dec}
    \mathbf{a} + \mathbf{b}_{\rm dec}$
\end{enumerate}

\noindent di mana $\mathbf{W}_{\rm enc} \in \mathbb{R}^{m \times d}$ (dengan
$m \gg d$) adalah matriks \emph{encoder}, $\text{TopK}(\cdot, k)$ adalah
fungsi yang mempertahankan hanya $k$ aktivasi terbesar dan menol-kan sisanya
(mendorong \emph{sparsity}), dan $\mathbf{W}_{\rm dec} \in \mathbb{R}^{d \times m}$
adalah matriks \emph{decoder} dengan kolom ternormalisasi.

Fungsi kerugian SAE didefinisikan sebagai:

\begin{equation}\label{eq:bab3_sae_loss}
  \mathcal{L}_{\rm SAE}(\theta) = \|\mathbf{z} - \hat{\mathbf{z}}\|_2^2,
\end{equation}

\noindent yaitu rekonstruksi MSE tanpa penalti $L_1$ eksplisit (karena
\emph{sparsity} sudah dipaksakan melalui TopK). Setelah pelatihan, setiap
kolom $\mathbf{w}_j^{\rm dec}$ dari $\mathbf{W}_{\rm dec}$ merupakan vektor
fitur laten yang dapat diinterpretasi, dan nilai $a_j$ mengindikasikan
seberapa kuat fitur ke-$j$ diaktivasi untuk input tertentu
\citetitb{Bricken2023Monosemanticity, Templeton2024Scaling}.

Gambar~\ref{fig:bab3_sae_topology} mengilustrasikan topologi SAE dengan
aktivasi Top-$k$. Bottleneck ekspansi (dimensi laten $m \gg d$) mendorong
pemisahan konsep, sementara fungsi TopK memastikan hanya $k$ dari $m$ unit
yang aktif pada setiap langkah waktu, menjamin \emph{monosemanticity}
fitur yang terpelajar.

\begin{figure}[ht]
  \centering
  \resizebox{0.9\textwidth}{!}{%
  \begin{tikzpicture}[
      node distance = 0.45cm and 2.4cm,
      unit/.style  = {circle, draw, thick, fill=white, minimum size=0.5cm},
      active/.style = {unit, fill=blue!25},
      box/.style   = {rectangle, draw, thick, fill=gray!10, rounded corners=3pt,
                      text centered, minimum height=0.6cm, font=\small},
      arr/.style   = {-{Stealth[length=4pt]}, thick},
      sarr/.style  = {arr, gray!70},
    ]
    \foreach \i in {1,...,5}{
      \node[unit] (in\i) at (0, {-(\i-1)*0.7}) {};
    }
    \node[box, left=0.6cm of in3] (zinput) {$\mathbf{z}\in\mathbb{R}^d$};

    \foreach \j in {1,2,4,5,7,8}{
      \node[unit]   (lat\j) at (3.0, {-(\j-1)*0.5+0.35}) {};
    }
    \node[active] (lat3) at (3.0, {-(3-1)*0.5+0.35}) {};
    \node[active] (lat6) at (3.0, {-(6-1)*0.5+0.35}) {};
    \node[box, above=0.2cm of lat1, text=blue!60!black, font=\scriptsize]
         {Top-$k$ aktif};

    \foreach \i in {1,...,5}{
      \node[unit] (out\i) at (6.0, {-(\i-1)*0.7}) {};
    }
    \node[box, right=0.6cm of out3] (zrec) {$\hat{\mathbf{z}}\in\mathbb{R}^d$};

    \foreach \i in {1,...,5}{
      \draw[sarr] (in\i.east) -- (lat3.west);
      \draw[sarr] (in\i.east) -- (lat6.west);
    }
    \foreach \i in {1,...,5}{
      \draw[sarr] (lat3.east) -- (out\i.west);
      \draw[sarr] (lat6.east) -- (out\i.west);
    }

    \node[font=\small, below=0.3cm of in5]  {Encoder $\mathbf{W}_{\rm enc}$};
    \node[font=\small, below=0.3cm of out5] {Decoder $\mathbf{W}_{\rm dec}$};
    \node[font=\scriptsize, below=0.35cm of lat8] {$m \gg d$};
  \end{tikzpicture}%
  }
  \caption{Topologi \emph{Top-$k$ Sparse Autoencoder} (SAE). Representasi laten
           $\mathbf{z} \in \mathbb{R}^d$ dari model RUL diencode ke ruang yang
           lebih besar ($m \gg d$), kemudian fungsi TopK mempertahankan hanya $k$
           unit teraktif (biru) dan menolkan sisanya. Matriks decoder merekonstruksi
           $\hat{\mathbf{z}}$, dan setiap kolom $\mathbf{w}_j^{\rm dec}$ merupakan
           satu \emph{fitur laten} yang dapat diinterpretasi.}
  \label{fig:bab3_sae_topology}
\end{figure}

\subsection{SHAP (\emph{SHapley Additive exPlanations})}
\label{subsec:bab3_shap}

Nilai SHAP untuk fitur $i$ pada prediksi $f(\mathbf{x})$ didefinisikan
sebagai nilai Shapley:

\begin{equation}\label{eq:bab3_shap}
  \phi_i(f, \mathbf{x}) = \sum_{S \subseteq \mathcal{F} \setminus \{i\}}
    \frac{|S|!(|\mathcal{F}| - |S| - 1)!}{|\mathcal{F}|!}
    \left[ v(S \cup \{i\}) - v(S) \right],
\end{equation}

\noindent di mana $\mathcal{F}$ adalah himpunan semua fitur, $S$ adalah
subset fitur yang mungkin, dan $v(S) = \mathbb{E}[f(\mathbf{x}) \mid
\mathbf{x}_S]$ adalah prediksi yang diharapkan ketika hanya subset $S$ yang
diketahui \citetitb{Lundberg2017SHAP}. Nilai SHAP memenuhi sifat-sifat
penjumlahan, tidak hadir, simetri, dan dumminess yang diinginkan secara
aksiomatik.

\subsection{\emph{Integrated Gradients} (IG)}
\label{subsec:bab3_ig}

\emph{Integrated Gradients} mendefinisikan atribusi fitur $i$ sebagai:

\begin{equation}\label{eq:bab3_ig}
  \text{IG}_i(\mathbf{x}) = (x_i - x'_i)
    \int_{0}^{1}
      \frac{\partial f\!\left(\mathbf{x}' + \alpha(\mathbf{x} - \mathbf{x}')\right)}
           {\partial x_i}
    \, \mathrm{d}\alpha,
\end{equation}

\noindent di mana $\mathbf{x}'$ adalah titik \emph{baseline} (biasanya vektor
nol) dan $\alpha$ adalah parameter interpolasi \citetitb{Sundararajan2017IG}.
Integral ini diestimasi menggunakan kuadratur trapesoid dengan $m$ langkah
interpolasi. IG memenuhi dua aksiom kunci: \emph{sensitivity} (jika model
bergantung pada fitur $i$, atribusi $i$ tidak nol) dan
\emph{implementation invariance} (arsitektur yang secara fungsional setara
menghasilkan atribusi yang sama).

% -----------------------------------------------------------------
\section{Metrik Evaluasi Prognostik}
\label{sec:bab3_metrik}
% -----------------------------------------------------------------

Evaluasi kinerja model prediksi RUL menggunakan empat metrik yang masing-masing
menyorot aspek berbeda dari akurasi prediksi. Metrik-metrik ini dihitung pada
set data uji setelah model selesai dilatih.

\subsection{\emph{Root Mean Squared Error} (RMSE)}
\label{subsec:bab3_rmse}

RMSE mengukur deviasi kuadratik rata-rata antara prediksi dan nilai sebenarnya:

\begin{equation}\label{eq:bab3_rmse}
  \text{RMSE} = \sqrt{\frac{1}{N} \sum_{i=1}^{N} (r_i - \hat{r}_i)^2},
\end{equation}

\noindent di mana $r_i$ dan $\hat{r}_i$ masing-masing adalah target dan
prediksi RUL pada sampel ke-$i$, dan $N$ adalah jumlah total sampel.
RMSE sensitif terhadap kesalahan besar dan merupakan metrik primer dalam
perbandingan antar model.

\subsection{\emph{Mean Absolute Error} (MAE)}
\label{subsec:bab3_mae}

MAE mengukur deviasi absolut rata-rata dan lebih \emph{robust} terhadap
\emph{outlier}:

\begin{equation}\label{eq:bab3_mae}
  \text{MAE} = \frac{1}{N} \sum_{i=1}^{N} |r_i - \hat{r}_i|.
\end{equation}

\subsection{Koefisien Determinasi ($R^2$)}
\label{subsec:bab3_r2}

$R^2$ mengukur fraksi varians target yang dapat dijelaskan oleh model:

\begin{equation}\label{eq:bab3_r2}
  R^2 = 1 - \frac{\sum_{i=1}^{N} (r_i - \hat{r}_i)^2}{\sum_{i=1}^{N} (r_i - \bar{r})^2},
\end{equation}

\noindent di mana $\bar{r}$ adalah nilai rata-rata target. Perlu dicatat bahwa
pada skema pelabelan \emph{piecewise} di mana mayoritas nilai target bernilai 1
(fase sehat panjang), varians target kecil sehingga $R^2$ menjadi kurang
informatif; RMSE dan PHM Score lebih diprioritaskan dalam analisis.

\subsection{PHM Score (Fungsi Skor Asimetris)}
\label{subsec:bab3_phm_score}

Fungsi skor standar kompetisi IEEE PHM 2012 \citetitb{Nectoux2012PRONOSTIA}
memberikan penalti asimetris: prediksi yang terlalu dini (optimistik)
mendapat penalti lebih ringan dibandingkan prediksi yang terlambat
(pesimistik), karena secara operasional lebih aman memperkirakan kegagalan
lebih cepat:

\begin{align}
  E_i &= 100 \cdot \frac{r_i - \hat{r}_i}{r_i}, \label{eq:bab3_phm_ei} \\
  A_i &= \begin{cases}
    \exp\!\left(-\ln(0{,}5) \cdot \dfrac{E_i}{5}\right), & E_i \leq 0 \\[6pt]
    \exp\!\left(+\ln(0{,}5) \cdot \dfrac{E_i}{20}\right), & E_i > 0
  \end{cases}, \label{eq:bab3_phm_ai} \\
  \text{PHM Score} &= \frac{1}{N} \sum_{i=1}^{N} A_i. \label{eq:bab3_phm_score}
\end{align}

Nilai PHM Score maksimum adalah 1,0 (prediksi sempurna). Prediksi yang
terlambat ($E_i < 0$, model mengestimasi RUL terlalu tinggi dibanding
aktualnya) dihukum lebih berat karena dapat menyebabkan operasi \emph{bearing}
melewati titik kegagalan tanpa intervensi. Interpretasi lengkap dan analisis
distribusi PHM Score per \emph{bearing} disajikan pada
Bab~\ref{bab:hasil-pembahasan}.

% -----------------------------------------------------------------
\section{Klasifikasi Kondisi \emph{Bearing} dan Algoritma Kandidat (Pilar~0)}
\label{sec:bab3_klasifikasi}
% -----------------------------------------------------------------

Subbab ini memuat fondasi teoretis untuk Pilar~0 disertasi --- studi
kasus industri yang melakukan klasifikasi tiga kelas kondisi \emph{bearing}
(\emph{Normal}, \emph{Warning}, \emph{Critical}) pada lini produksi
PT~SKF~Indonesia. Dasar teori ini diletakkan terpisah dari materi RUL
karena keluaran model, fungsi kerugian, dan metrik evaluasi berbeda
secara struktural.

\subsection{Formulasi Klasifikasi Multi-Kelas}
\label{subsec:bab3_kls_formulasi}

Diberikan jendela sekuens fitur $\mathbf{X} \in \mathbb{R}^{L \times d}$
($L$ panjang jendela, $d$ jumlah fitur per langkah), pengklasifikasi
$f_\theta : \mathbb{R}^{L \times d} \to \Delta^{K-1}$ memetakan jendela
ke distribusi probabilitas atas $K$ kelas pada simpleks
$\Delta^{K-1}$. Untuk Pilar~0, $K = 3$ dengan label
$\{\,\text{Normal},\, \text{Warning},\, \text{Critical}\,\}$. Prediksi
kelas diperoleh dengan
$\hat{y} = \arg\max_{k \in \{1,\ldots,K\}} f_\theta(\mathbf{X})_k$, dan
parameter dilatih dengan meminimalkan \emph{cross-entropy}:

\begin{equation}\label{eq:bab3_xentropy}
  \mathcal{L}_{\text{CE}}(\theta) = -\frac{1}{N} \sum_{i=1}^{N}
    \sum_{k=1}^{K} y_{i,k} \log f_\theta(\mathbf{X}_i)_k,
\end{equation}

\noindent dengan $y_{i,k} \in \{0,1\}$ adalah label \emph{one-hot} sampel
ke-$i$. Karena distribusi kelas pada data industri umumnya tidak
seimbang (kelas \emph{Critical} jarang teramati), digunakan
\emph{class-weighted cross-entropy} dengan bobot
$w_k \propto N / (K \cdot N_k)$.

\subsection{Metrik Evaluasi Klasifikasi}
\label{subsec:bab3_kls_metrik}

Kinerja klasifikasi dievaluasi dengan empat metrik yang saling melengkapi:

\begin{equation}\label{eq:bab3_kls_metrik}
  \begin{aligned}
    \text{Accuracy} &= \frac{\text{TP} + \text{TN}}{\text{TP} + \text{TN} + \text{FP} + \text{FN}}, \\
    \text{Precision}_k &= \frac{\text{TP}_k}{\text{TP}_k + \text{FP}_k},\quad
    \text{Recall}_k = \frac{\text{TP}_k}{\text{TP}_k + \text{FN}_k}, \\
    F1_k &= 2 \cdot \frac{\text{Precision}_k \cdot \text{Recall}_k}{\text{Precision}_k + \text{Recall}_k},
  \end{aligned}
\end{equation}

\noindent dengan $F1_{\text{macro}} = (1/K)\sum_k F1_k$ menjadi metrik
utama karena tidak terbias terhadap kelas mayoritas. Untuk konteks
pemeliharaan, \emph{Recall} kelas \emph{Critical} adalah metrik
keselamatan paling kritis: kegagalan mendeteksi kelas ini berakibat
\emph{unplanned downtime}. Dua metrik operasional juga dilaporkan:
\emph{inference latency} per jendela ($\leq 10$~ms untuk implementasi
\emph{near real-time}) dan jumlah parameter terlatih sebagai proksi
\emph{memory footprint}.

\subsection{Algoritma Kandidat untuk Benchmarking}
\label{subsec:bab3_kls_kandidat}

Enam algoritma yang dibandingkan pada Pilar~0 dipilih untuk
merepresentasikan tiga keluarga pendekatan: ML tradisional,
\emph{deep learning} murni, dan hibrida. Mekanisme detail tiap
algoritma sudah dibahas sebagian pada bab ini (LSTM, GRU, mLSTM,
SSM); ringkasan komplementer per kandidat dijelaskan di bawah.

\paragraph{XGBoost (\emph{baseline} ML tradisional).}
\emph{Extreme Gradient Boosting} \citetitb{ChenGuestrin2016XGBoost}
adalah \emph{ensemble} pohon keputusan yang dilatih secara
\emph{additive}: pada iterasi ke-$m$ ditambahkan satu pohon $h_m$ yang
meminimalkan
\begin{equation*}
  \mathcal{L}^{(m)} = \sum_i \ell\bigl(y_i,\, \hat{y}_i^{(m-1)} + h_m(\mathbf{x}_i)\bigr)
    + \Omega(h_m),
\end{equation*}
dengan $\Omega(h) = \gamma T + \tfrac{1}{2}\lambda\|w\|^2$ sebagai
regularisasi pada jumlah daun $T$ dan bobot daun $w$. Input model adalah
vektor fitur \emph{handcrafted} (RMS, kurtosis, \emph{crest factor},
energi pita BPFx, dan fitur kualitas produk).

\paragraph{CNN-LSTM \emph{end-to-end}.}
Dua-tahap: lapisan konvolusi 1D mengekstraksi fitur spasial dari
segmen sinyal getaran mentah, lalu LSTM menangkap dependensi temporal
\citetitb{Khorram2020CNNLSTM}. Mekanisme LSTM dijelaskan pada
Persamaan~\ref{eq:bab3_lstm_i}--\ref{eq:bab3_lstm_h}. Lapisan akhir adalah
\emph{fully-connected} dengan \emph{softmax} $K$ kelas.

\paragraph{Temporal Convolutional Network (TCN).}
TCN \citetitb{Bai2018TCN} memakai \emph{dilated causal convolution}
\begin{equation*}
  y_t = \sum_{k=0}^{K-1} f(k)\, x_{t - d \cdot k},
\end{equation*}
dengan faktor dilatasi $d \in \{1, 2, 4, \ldots\}$ tumbuh eksponensial
sehingga \emph{receptive field} mencakup beberapa ribu langkah waktu
tanpa beban rekuren. \emph{Residual block} ditambahkan agar pelatihan
stabil pada kedalaman besar.

\paragraph{BiLSTM dengan \emph{Hierarchical Improved Fusion Attention}.}
\citetitb{Huang2025BiLSTMHIF} menggabungkan BiLSTM dengan dua mekanisme
\emph{attention}: \emph{Time-Frequency Improved Attention} (T-FIAttn)
yang membobotkan fitur waktu-frekuensi, dan \emph{Layered Global
Adaptive Attention Mechanism} (L-GAAM) yang membobotkan fitur antar
lapisan. \emph{Gating} di akhir memungkinkan BiLSTM fokus pada subset
fitur yang paling esensial untuk klasifikasi.

\paragraph{GRU ringan untuk \emph{edge}.}
\emph{Gated Recurrent Unit} \citetitb{Cho2014GRU} menggabungkan
\emph{forget} dan \emph{input gate} LSTM menjadi satu \emph{update
gate}, sehingga jumlah parameter berkurang $\approx 33\%$. Pada Pilar~0,
GRU dipasangkan dengan kanal ganda 1D-CNN paralel mengikuti
\citetitb{Bai2025ParallelGRU} agar tetap kompetitif tanpa kehilangan
keuntungan ringan untuk \emph{deployment} di SKF IMx-8 atau perangkat
\emph{edge} setara.

\paragraph{Hibrida ARIMA-LSTM.}
Model \emph{Autoregressive Integrated Moving Average} (ARIMA)
\citetitb{Box1976ARIMA} memodelkan komponen linier deret waktu:
\begin{equation*}
  \phi(B)\, (1-B)^d\, x_t = \theta(B)\, \varepsilon_t,
\end{equation*}
dengan $B$ \emph{backshift operator} dan $\phi, \theta$ polinomial
\emph{autoregressive} dan \emph{moving-average}. Residu non-linier
$r_t = x_t - \hat{x}_t^{\text{ARIMA}}$ diumpankan ke LSTM untuk
menangkap pola yang tidak dapat dipotret ARIMA \citetitb{Hamiane2024ARIMALSTM}.

\subsection{SHAP untuk Klasifikasi: Implementasi Praktis}
\label{subsec:bab3_kls_shap}

SHAP sudah diuraikan pada subbab~\ref{subsec:bab3_shap}. Dalam konteks
Pilar~0, dua varian dipakai: \emph{TreeExplainer} untuk XGBoost
(eksak, $O(TLD^2)$ untuk $T$ pohon, $L$ daun, $D$ kedalaman maksimum)
dan \emph{KernelExplainer}/\emph{DeepExplainer} untuk DL. Output utama
adalah peringkat \emph{global feature importance}
$\bar{\phi}_j = \tfrac{1}{N}\sum_i |\phi_{i,j}|$ dan plot
\emph{summary} per kelas. Penjelasan per-prediksi (\emph{waterfall plot})
dipakai untuk \emph{root-cause analysis} pada misklasifikasi
\emph{Critical} yang berdampak operasional.

\subsection{Skema Pelabelan Tiga Kelas}
\label{subsec:bab3_kls_label}

Pelabelan kondisi tiga kelas mengikuti praktik standar SKF Observer
yang membagi rentang \emph{Progression Index} (PI) ternormalisasi
ke ambang berikut: \emph{Normal} jika $\text{PI} < 0{,}20$; \emph{Warning}
jika $0{,}20 \leq \text{PI} < 0{,}50$; dan \emph{Critical} jika
$\text{PI} \geq 0{,}50$. PI sendiri merupakan rasio kenaikan
parameter getaran dari \emph{baseline} yang dikoreksi oleh kemiringan
tren \emph{moving window}. Konsistensi label tervalidasi melalui
\emph{cross-check} dengan ambang ISO~10816/20816 untuk kelas
\emph{velocity RMS} mesin grinding industri presisi.
